\section{Conclusion}\label{sec:conclusion}

We fitted polynomials of degree up to 20 to noisy samples of Runge's function
with OLS, Ridge and Lasso, selected models with the bootstrap and $k$-fold
cross-validation, and solved the regression problems with our own
gradient-descent code. The central finding is that the resampling methods, not
the regression methods, limit model selection here. The true test error of OLS
is flat over a wide range of high degrees, but every resampled fit must
extrapolate to held-out points outside its training range, where a high-degree
polynomial diverges, a statistical echo of Runge's phenomenon. The estimators
therefore select too low a degree, the more so the fewer distinct points each
fit uses. Ridge suppresses the directions of small singular values that cause
these failures, which lets CV choose higher degrees, but the selected penalties
are tiny and the gain on held-out data is negligible. Lasso did worse than OLS
because our subgradient solver does not reach the Lasso solution at high
degree. On the optimization side, the condition number of the Hessian, which
grows exponentially with the degree, explains our results: plain GD needs about
$\kappa$ iterations, momentum and Adam are an order of magnitude faster, no
first-order method reaches the OLS solution at $p\geq10$, a small Ridge penalty
restores convergence, and subgradient descent with a fixed step levels off at a
distance from the optimum set by the step.

Our study has clear limitations. The true error, which we used to judge the
estimators, requires the known $f$ and is unavailable in practice; the same
holds for the bias against $f$. We studied a single noise level and sample
size for the model selection, a uniform random design, and one set of folds per
data set, so the size of the extrapolation effect is specific to these
choices. The bootstrap estimates of bias and
variance are means over replicates and are dominated by the few replicates that
extrapolate badly; a robust estimate (median) would separate the two effects.
The optimizer comparison used a learning-rate grid with four values per decade
and fixed momentum and Adam parameters, which is too coarse to rank momentum
and Adam where they differ by less than a factor of two, and it measures
iterations, not time; for SGD we used a single learning-rate decay. Our Lasso results are limited by the solver: with a fixed
number of 5000 Adam steps, the high-degree Lasso fits in the cross-validation
are far from converged, so the Lasso model selection and comparison measure our
optimizer as much as the Lasso model itself. Finally,
all comparisons on the held-out data rest on 20 test points, whose MSE has a
standard error of about $30\%$ of $\sigma^2$.

Two extensions follow directly from these limitations. First, replacing the
monomials $x^j$ by orthogonal polynomials on $[-1,1]$ spans the same model
space but keeps $\kappa(\bm{X})$ small; gradient descent, including the
subgradient method for Lasso, would then converge at every degree, and the
coefficients would become interpretable. Second, the resampling estimates
could be made less sensitive to extrapolation by stratifying the folds in $x$,
so that every training set spans the whole interval, or by using leave-one-out
cross-validation, which for OLS follows from a single fit. Whether the selected
degree then approaches the true optimum would show how much of the gap in
\cref{fig:ols_estimators} is due to extrapolation.
